\documentclass[11pt,letterpaper]{article}
\usepackage[T1]{fontenc}
\usepackage[utf8]{inputenc}
\usepackage{lmodern,microtype}
\usepackage[margin=1in,headheight=14pt]{geometry}
\usepackage{amsmath,amssymb,amsthm,mathtools}
\usepackage{booktabs,array,longtable,enumitem}
\usepackage[dvipsnames]{xcolor}
\usepackage{hyperref}
\usepackage{fancyhdr}
\hypersetup{colorlinks=true,linkcolor=MidnightBlue,citecolor=MidnightBlue,urlcolor=MidnightBlue,
 pdftitle={Sharp Moment Reconstruction at Surreal Scales},
 pdfauthor={Research manuscript prepared with ChatGPT}}
\pagestyle{fancy}\fancyhf{}
\fancyhead[L]{\small Moment reconstruction at surreal scales}
\fancyhead[R]{\small Research manuscript}
\fancyfoot[C]{\thepage}
\setlength{\parskip}{4pt}
\setlength{\parindent}{1.2em}
\setlist{nosep,leftmargin=*}
\allowdisplaybreaks[2]
\newtheorem{theorem}{Theorem}[section]
\newtheorem{lemma}[theorem]{Lemma}
\newtheorem{proposition}[theorem]{Proposition}
\newtheorem{corollary}[theorem]{Corollary}
\theoremstyle{definition}
\newtheorem{definition}[theorem]{Definition}
\newtheorem{example}[theorem]{Example}
\theoremstyle{remark}
\newtheorem{remark}[theorem]{Remark}
\newcommand{\C}{\mathbb C}
\newcommand{\R}{\mathbb R}
\newcommand{\Q}{\mathbb Q}
\newcommand{\N}{\mathbb N}
\newcommand{\Z}{\mathbb Z}
\newcommand{\No}{\mathbf{No}}
\newcommand{\SC}{\mathbf{No}[i]}
\newcommand{\Ocal}{\mathcal O}
\newcommand{\mfr}{\mathfrak m}
\newcommand{\supp}{\operatorname{supp}}
\newcommand{\lc}{\operatorname{lc}}
\newcommand{\st}{\operatorname{st}}
\newcommand{\vg}{v_{\mathrm G}}
\newcommand{\Ph}{\widehat P}
\newcommand{\Ah}{\widehat A}
\newcommand{\ah}{\widehat a}
\newcommand{\wh}{\widehat w}
\newcommand{\mh}{\widehat m}
\newcommand{\Lh}{\widehat L}
\newcommand{\Id}{\mathrm I}
\newcommand{\mom}{\mathcal M}
\newcommand{\dd}{\mathrm d}
\newcommand{\lbl}[1]{\nolinkurl{#1}}% label names and paths of the collection, breakable at : and /
\DeclareMathOperator{\diag}{diag}
\DeclareMathOperator{\Disc}{Disc}
\DeclareMathOperator{\Res}{Res}
\DeclareMathOperator{\ord}{ord}
\DeclareMathOperator{\rank}{rank}
\title{\textbf{Sharp Moment Reconstruction\\at Surreal Scales}\\[5pt]
\large Collision geometry, optimal precision, and Hahn-supported inversion}
\author{Research manuscript prepared with ChatGPT}
\date{22 September 2026}
\begin{document}
\maketitle
\begin{abstract}
We study recovery of a finite weighted configuration from its first $2n$ power
moments over real and complex Hahn fields, and hence over the surreal and
surcomplex numbers. For distinct finite nodes $a_i$ and nonzero weights $w_i$,
put $d_i=\sum_{j\ne i}v(a_i-a_j)$ and
$\delta_i=\max_{j\ne i}v(a_i-a_j)$. We prove that
\[
 \Theta=\max_i\bigl(v(w_i)+2d_i+\delta_i\bigr)
\]
is the exact uniform valuation threshold for guaranteeing reconstruction
inside all strict nearest-neighbour balls. Errors of valuation at least
$\kappa>\Theta$ admit a unique regular $n$-node realization, with node error
at least $\kappa-v(w_i)-2d_i$ and weight error at least
$\kappa-2d_i-\delta_i$. Perturbing only the last moment attains every node bound
simultaneously; at valuation $\Theta$ it defeats the full strict-ball
guarantee for every starting configuration. We also give an arbitrary-rank
precision-lattice theorem, quantify cancellation in the inverse Hermite
matrix, and derive a directed-path certificate for nonuniform moment
precisions. The arguments use finite algebra and strong Hahn summability,
not an Archimedean contraction argument. The directed-path certificate also
holds for nondivisible value groups: fractional scaling exponents are
auxiliary, and the reconstruction stays in the original Hahn field. Classical Prony reconstruction,
Hermite interpolation, and differential precision tracking are explicitly
separated from the proposed refinements. Priority is not certified; the
article supplies proofs rather than claiming a solution to a named
published conjecture.

In this collection the report is built from a single manuscript and filed
under Prony's name, because ``moment'' has other meanings in neighbouring
reports. Section~\ref{prony:subsec:collection} records its provenance, fixes
the meaning of ``moment'' used here, and compares $\Theta$ with the
coefficient threshold $T(P)=\max_i(d_i+\delta_i)$ of the collection's
polynomial-algebra report.
\end{abstract}
\tableofcontents
\newpage

\section{The problem, the contribution, and its scope}
\label{prony:sec:intro}

Fix an integer $n\ge1$. A finite atomic configuration is specified by distinct nodes $a_1,\ldots,a_n$
and nonzero weights $w_1,\ldots,w_n$. Its moments are
\begin{equation}\label{prony:eq:moments}
 m_k=\sum_{i=1}^n w_i a_i^k,\qquad 0\le k<2n.
\end{equation}
Section~\ref{prony:subsec:collection} fixes this meaning of ``moment'' for the
whole report and compares it with other uses of the word in the collection.
For ordinary complex data this is the classical Prony problem. When two
nodes are extremely close, inversion magnifies measurement errors. Over
surreal scalars, closeness can occur at several incomparable Archimedean
orders: one cluster may be separated at scale $\omega^{-1}$, another at
$\omega^{-\omega}$, and nested clusters may combine both. A single real
condition number cannot retain that ordered-group information.

The question addressed here is precise: \emph{how much moment precision
forces the reconstructed nodes to retain their individual labels and their
entire finite separation geometry?} We do not merely seek continuity or a
first-order condition number. We seek a sufficient threshold, exact node
losses above it, and a counterexample at the threshold itself.

Throughout, ``regular $n$-node realization'' means exactly $n$ distinct nodes
with nonzero weights. Uniqueness is in this fixed-order model; $2n$ moments
do not determine arbitrary configurations with more than $n$ nodes.

\subsection{Statement of the main result}
Let $K=\Bbbk((t^\Gamma))$, where $\Bbbk=\R$ or $\C$ and $\Gamma$ is a
nontrivial ordered abelian group. The support of a Hahn series is
well-ordered, and $v(t^\gamma)=\gamma$. Write
$\Ocal=\{x:v(x)\ge0\}$ and $\mfr=\{x:v(x)>0\}$, with $v(0)=+\infty$.
Assume $a_i\in\Ocal$, and define
\begin{equation}\label{prony:eq:geometry}
\begin{aligned}
 d_{ij}&=v(a_i-a_j),&
 d_i&=\sum_{j\ne i}d_{ij},& \delta_i&=\max_{j\ne i}d_{ij},\\
 \varpi_i&=v(w_i),& E_i&=\varpi_i+2d_i,&
 \Theta&=\max_i(E_i+\delta_i).
\end{aligned}
\end{equation}
For $n=1$ the empty sum and the empty maximum in this notation are both
assigned the value zero. The convention fixes the residue ball in that case.
The letters $d_{ij}$, $d_i$ and $\delta_i$ have exactly the meaning they have
in the collection's polynomial-algebra report; the weight valuations are
written $\varpi_i$, because $\alpha_i$ denotes roots there
(Section~\ref{prony:subsec:collection}).

\begin{theorem}[Optimal uniform reconstruction]\label{prony:thm:main}
Suppose $\mh_k=m_k+\epsilon_k$, with $v(\epsilon_k)\ge\kappa>\Theta$ for
$0\le k<2n$. There is a unique regular $n$-node realization of the perturbed
moments, up to permutation. It has a unique labelling satisfying
\begin{equation}\label{prony:eq:strictballs}
 v(\ah_i-a_i)>\delta_i.
\end{equation}
With this labelling,
\begin{align}
 v(\ah_i-a_i)&\ge\kappa-E_i,\label{prony:eq:nodebound}\\
 v(\wh_i-w_i)&\ge\kappa-2d_i-\delta_i>\varpi_i.\label{prony:eq:weightbound}
\end{align}
Consequently the leading terms of the weights and of every nonzero
pairwise node difference are preserved.

If only $m_{2n-1}$ is perturbed, by a nonzero $e$ of valuation
$\kappa>\Theta$, equality holds in \eqref{prony:eq:nodebound} for every $i$.
If instead $v(e)=\Theta$, no regular $n$-node realization of those moments
can satisfy \eqref{prony:eq:strictballs} for all $i$.
\end{theorem}

The final assertion means that the open precision condition is genuinely
optimal for \emph{this} labelling guarantee. It does not assert that every
boundary perturbation causes a multiple root, or that reconstruction is
impossible below the threshold. Some perturbations have much more
cancellation than the worst case.

\subsection{Relation to the literature}
Prony reconstruction by a Hankel system, polynomial roots, and a Vandermonde
system is classical. Batenkov and Yomdin study local accuracy of confluent
Prony systems \cite{BY}. Katz, Diab, and Batenkov obtain sharp clustered
complex reconstruction estimates and analyze correlations between the
successive errors \cite{KDB}. In particular, their first-order node formula
already has a denominator consisting of the weight times squared pairwise
separations. Thus the occurrence of $w_iP'(a_i)^2$ below is not presented as
a previously unknown conditioning mechanism. For a single cluster of $n$
unit-weight nodes with all separation valuations equal to $s$, our formulas
give the familiar node loss $(2n-2)s$ and threshold $(2n-1)s$.

The refinement proposed here is the \emph{exact ordered-group statement}
for arbitrary finite nested separation patterns, weights at arbitrary
valuations, and an optimality witness valid at every configuration. The
threshold is pointwise rather than a cluster-size asymptotic. It applies
without embedding the value group into $\R$ and without restricting nodes
to an ordinary unit circle. The proof isolates the correlations by changing
to the cofactor basis before making any estimates.


There is also a classical geometric antecedent to the sharpness direction:
Akinshin--Goldman--Yomdin study Prony varieties obtained by fixing initial
moments, including the curve on which the first $2n-1$ moments stay fixed
\cite{AGY}. Thus varying only the last measured moment is not claimed as
a newly discovered direction. The proposed refinement is its exact
valuation consequence for every arbitrary-rank configuration, together
with the uniform converse.

Exact propagation of precision lattices by a differential is an established
non-Archimedean method, notably in Caruso--Roe--Vaccon \cite{CRV}. Our
Section~\ref{prony:sec:lattice} is a finite-polynomial, arbitrary-rank version
proved by Hahn support evaluation. Neither its general philosophy nor
Hermite interpolation is claimed as new. Its use here provides exact
nonlinear uncertainty sets and exposes a cancellation invariant not
recorded by the separation tree alone. The nonuniform certificate of
Section~\ref{prony:sec:graph} combines the cofactor system with directed-path
bounds; it is a sufficient certificate, not an exact characterization.

\subsection{Relation to the supplied repository}
\label{prony:subsec:repository}
The comparison uses the repository snapshot
\begin{center}
\texttt{4cf691c7d951e037739d32d9f5c387dcce724f3c}.
\end{center}
The documentation catalogue, formalization-ledger excerpts, and the
polynomial-algebra report's detailed README were inspected \cite{Repo}.
That report already contains finite Hermite jets, Hahn factor lifting,
root-cluster geometry, and a \emph{coefficient-perturbation} threshold of
the form $\max_i(d_i+\delta_i)$. (In that report's current labels these are
\lbl{polynomial:thm:crt}, the section \lbl{polynomial:sec:hensel}, the
cluster tree of \lbl{polynomial:subsec:tree}, and
\lbl{polynomial:thm:stability} \cite{RepoPoly}.) Those are antecedents, not
contributions of this
article. The present input is a vector of \emph{moments}, whose structured
inversion changes the threshold to
$\max_i(\varpi_i+2d_i+\delta_i)$ and requires control of the recovered weights.
Remark~\ref{prony:rem:coefficient} states the relation between the two
thresholds precisely.

At the pinned snapshot, repository searches for ``Prony'' and ``Hankel''
returned no indexed matches, and the separate \lbl{docs/new} tree listed
eighteen ZIP archives; their
binary contents and all historical revisions were not exhaustively read.
Accordingly, ``not found in the inspected repository material'' is the
scope of the nonduplication check. Both observations are provenance, not
current descriptions. The first remained true of the collection's reports
until this batch was placed, but the collection now also contains the
Hankel and Prony-type material compared in
Section~\ref{prony:subsec:collection}. The archives that were then in
\lbl{docs/new}, and the later ones including the archive that delivered this
manuscript, have since been placed into reports and removed from that
directory. The literature comparison is likewise
not a proof of priority. The strongest candidate contribution is
Theorem~\ref{prony:thm:main} together with its universal last-moment sharpness
mechanism; the other statements are useful supporting refinements.

\subsection{This report in the collection}
\label{prony:subsec:collection}

\paragraph{Provenance.}
This report is built from one manuscript: the article ``Sharp Moment
Reconstruction at Surreal Scales'', dated 22 September 2026, with its
verification program and records. It was manuscript~13 of the eighteen placed
in the collection's eighteenth batch, delivered as the archive
\texttt{surcomplex\_moment\_reconstruction}, and its repository comparison is
pinned to commit \texttt{4cf691c} (the snapshot above). It is not a merge: no
other manuscript proves its main results, and the nearest material in the
collection, including a classical Prony-type lemma, is compared below. On placement the prefix \lbl{prony:} was
added to every label, and the four symbols in the table below were renamed.
This subsection, Section~\ref{prony:subsec:coefficient}, a note after
Proposition~\ref{prony:prop:prony}, pointers to other reports and two ledger
rows in Appendix~\ref{prony:app:ledger} were added. Statements about the
repository that were true only at the pin are now marked as such
(Sections~\ref{prony:subsec:repository} and~\ref{prony:sec:verification}
and Appendix~\ref{prony:app:audit}). The only choices were editorial: the
manuscript's title is kept, the directory is named after Prony, and the
notation follows the polynomial-algebra report. Subsequent proof corrections
are recorded in the maintained README. The limitations and priority caveats
are retained, and
the program \lbl{code/verify.py} and the files in \lbl{data/} are
unchanged from the delivery.
\begin{center}
\begin{tabular}{@{}llll@{}}
\toprule
Here & Manuscript & Meaning & In polynomial-algebra\\
\midrule
$d_{ij}$ & $s_{ij}$ & $v(a_i-a_j)$ & $d_{ij}$, the same\\
$\delta_i$ & $r_i$ & $\max_{j\ne i}d_{ij}$ & $\delta_i$, the same\\
$\varpi_i$ & $\alpha_i$ & $v(w_i)$ & not used; $\alpha_i$ are roots there\\
$\chi_i$ & $h_i$ & $\ell_i'(a_i)=\sum_{j\ne i}(a_i-a_j)^{-1}$ & not used\\
\bottomrule
\end{tabular}
\end{center}
The sum $d_i$ has the same meaning in both reports. The letter $\chi_i$ also
keeps the Hermite coefficient apart from $h=t^s$ in
Section~\ref{prony:sec:examples}. The program keeps the manuscript's letters
\texttt{alpha}, \texttt{r} and \texttt{hi}. Throughout, $\Theta$ is the
threshold \eqref{prony:eq:geometry}, an element of $\Gamma$; it has nothing to
do with the theta series of the collection's hahn-tate-uniformization
report. $T(P)$ always denotes the polynomial-algebra threshold; the letter $T$
inside the proofs of Proposition~\ref{prony:prop:lattice} and
Theorem~\ref{prony:thm:graph} is local.

\paragraph{What ``moment'' means here.}
In this report a moment is always one of the $2n$ power moments
\eqref{prony:eq:moments} of a finite weighted configuration, and the input
of reconstruction is that finite vector. The weights are arbitrary nonzero
elements of $K$; no measure and no positivity is involved, and
Corollary~\ref{prony:cor:positive} only records that positivity is preserved
when it is present. Two neighbouring reports use the word for other objects:
\lbl{surreal/hahn-valued-measures-and-probability} \cite{RepoMeas} for power
moments of Hahn-valued measures and moment functionals, including infinite
moment sequences and their representation problem, and
\lbl{surcomplex/hahn-herglotz-positivity} \cite{RepoHerg} for Fourier
(trigonometric) moments on the circle, tested by Toeplitz matrices. Power sums
of roots, that is, unit-weight moments with repeated roots allowed, are the
``contour moments'' of \lbl{d:thm:moments} in \lbl{surcomplex/analysis}
\cite{RepoAna} and the Newton sums of the Hankel section discussed below.
This is why the directory carries Prony's name rather than the word
``moment''.

\paragraph{The coefficient threshold.}
Theorem~\ref{prony:thm:main} is the moment-input analogue of
\lbl{polynomial:thm:stability} \cite{RepoPoly}, whose input is a
coefficient perturbation of a monic polynomial with distinct finite roots and
whose threshold is $T(P)=\max_i(d_i+\delta_i)$.
Remark~\ref{prony:rem:coefficient} shows that
$\Theta=\max_i\bigl((\varpi_i+d_i)+(d_i+\delta_i)\bigr)$, where
$\varpi_i+d_i=v(w_iP'(a_i))$ is the cost of passing from moments to the
annihilator, and that Lemma~\ref{prony:lem:localroots} is a nodewise form of
the matching step of that theorem.

\paragraph{Hankel matrices.}
With unit weights, the matrix $H$ of Proposition~\ref{prony:prop:prony} is
the Hankel matrix of the Newton power sums of $P$, that is, the matrix
$VV^{\mathsf T}$ of the trace form of $K[X]/(P)$, and $\det H=\Disc P$. This
classical identity is a common input to this report and the section on
Hankel square-class profiles in \lbl{surcomplex/spectral-theory}
\cite{RepoSpec}, which comes from another manuscript of the same batch and
states the identity as its \lbl{spec:hank:prop:hermite}. That
section reads square classes, modulo $2\Gamma$, of the leading principal
minors of such matrices over nondivisible value groups. This report
estimates valuation thresholds for perturbed weighted moments and uses
neither square classes nor leading principal minors. Their further theorems
concern different questions, and neither proof depends on the other report.

\paragraph{The measures report.}
\lbl{surreal/hahn-valued-measures-and-probability} \cite{RepoMeas}
contains, from another manuscript of the same batch, a Prony-type lemma on
finite signed exponential sums (a real sequence is a finite sum
$\sum_jc_jx_j^n$ with distinct real $x_j$ exactly when a monic polynomial
with distinct real roots annihilates it under the shift) and a finite Gaussian
quadrature theorem over real closed fields (a suitably positive linear
functional agrees up to degree $2N-1$ with $N$ positive weights at $N$
distinct nodes). Those results concern existence and exact representation of
moment data: the signed-sequence lemma requires no positivity, while
Gaussian quadrature does. This report starts from an $n$-node
configuration and studies the stability of its reconstruction, with exact
thresholds, and involves no measure. The results are cross-referenced, not
merged.

\section{Hahn fields, support, and surreal interpretation}
\label{prony:sec:hahn}

\subsection{Conventions and two different notions of smallness}
An element of $K$ is a formal sum
\[
 x=\sum_{\gamma\in S}x_\gamma t^\gamma,
 \qquad x_\gamma\in\Bbbk,
\]
with $S\subset\Gamma$ well-ordered. The valuation of a nonzero series is
its least exponent. Coefficient-zero extraction is a homomorphism
$\st:\Ocal\to\Bbbk$ with kernel $\mfr$. Every $x\in\Ocal$ has a unique
expression $x=x_0+x_+$ with $x_0\in\Bbbk$ and $x_+\in\mfr$.
For vectors, $v_{\min}(x)=\min_jv(x_j)$.

A family is \emph{strongly summable} when the union of its supports is
well-ordered and every exponent belongs to only finitely many supports.
The Hahn sum is then defined coefficientwise. This is not a claim that
finite partial sums converge in the full surreal fine topology. Nor does
$\eta>0$ imply that $N\eta$ eventually exceeds every element of $\Gamma$.
For example, in $\Q^2$ with lexicographic order, $N(0,1)<(1,0)$ for every
integer $N$. All infinite constructions used below are justified by support,
not by such a false cofinality assertion.

\begin{lemma}[Positive-support evaluation]\label{prony:lem:support}
If $z_1,\ldots,z_q\in\mfr$ and
$F\in\Bbbk[[Z_1,\ldots,Z_q]]$, its monomial expansion can be evaluated
strongly at $(z_1,\ldots,z_q)$. Evaluation preserves addition, multiplication,
and formal substitution when the substituted series have zero constant
term. The resulting support is contained in the additive monoid generated
by $\bigcup_j\supp(z_j)$.
\end{lemma}
\begin{proof}
Let $S$ be the finite union of the positive well-ordered supports, omitting
zero inputs. The word form of Higman's well-quasi-ordering theorem
\cite{Higman} implies the usual positive-support lemma: $S^*$, the monoid
of finite sums from $S$, is well-ordered, and a fixed exponent is the sum
of only finitely many finite words in $S$.
Here is the deduction. Order words by subsequence embedding with each letter
allowed to increase. An embedding never decreases the sum, and increases
it strictly unless the words are identical. An infinite decreasing sequence
of sums, or an infinite collection of distinct words with the same sum,
would therefore be a bad sequence for the word order. Higman's theorem
excludes both. The empty word supplies the sole sum of length zero.

For each input monomial, expand its product of Hahn series into words.
There are only finitely many choices of input labels for a fixed word.
The preceding finiteness is therefore exactly the local finiteness needed
to evaluate $F$. It also justifies the double-family regroupings in products
and zero-constant substitutions. No bound on the ordinary magnitudes of
the coefficients of $F$ is needed, because their valuations are zero.
\end{proof}

\begin{lemma}[Residue-simple lifting]\label{prony:lem:hensel}
If $f\in\Ocal[Y]$, $c\in\Bbbk$, and
$\overline f(c)=0$, $\overline f'(c)\ne0$, there is a unique
$y\in c+\mfr$ with $f(y)=0$. The same conclusion holds for a finite
polynomial system whose residue Jacobian at $c\in\Bbbk^N$ is invertible.
\end{lemma}
\begin{proof}
Write every coefficient as its residue plus an infinitesimal, and temporarily
replace the finitely many infinitesimal coefficients by independent formal
parameters $Z$. In the ring $\Bbbk[[Z]]$, solve for $y=c+u(Z)$ with $u(0)=0$
by total degree. At each positive degree the only contribution containing
the as-yet-unknown homogeneous part of $u$ is the invertible constant
Jacobian times that part. All other terms have already been determined.
This gives a unique formal solution. Evaluate it by
Lemma~\ref{prony:lem:support}.

For uniqueness in $c+\mfr^N$, subtract two equations. Polynomial divided
differences give a matrix multiplying their difference; its residue is the
invertible residue Jacobian. Its determinant is a unit of $\Ocal$, so the
difference is zero. This also proves uniqueness without an iterative-limit
argument.
\end{proof}

\subsection{Finite families in the full surreal field}
We use the standard Conway normal-form interpretation
\[
 \sum_\gamma c_\gamma t^\gamma
 \longmapsto \sum_\gamma c_\gamma\omega^{-\gamma}.
\]
For any set-sized ordered subgroup $\Gamma\subset\No$, real Hahn series
with exponents in $\Gamma$ identify with the corresponding normal forms;
complexification gives a subfield of $\SC$. The generalized-series
interpretation of surreal numbers is part of the standard background
reviewed in \cite{BG} and in the supplied repository.

Given finitely many actual surreal or surcomplex inputs, take the union of
their normal-form supports and the rational span of those exponents,
enlarging by $1$ when necessary. This is a set-sized divisible ordered
subgroup of $\No$. The associated Hahn field contains every input. Finite
field operations and the support evaluations above remain in that field.
If another workspace contains the same data, uniqueness of the constructed
solutions makes the two answers agree in a common enlargement.

\begin{corollary}[Actual surreal and surcomplex data]\label{prony:cor:actual}
Theorem~\ref{prony:thm:main}, the sharpness assertions, and all subsequent finite
precision statements apply to finite configurations in $\No$ or $\SC$,
using $v(\omega^{-\gamma})=\gamma$. Every construction is carried out in a
set-sized workspace. No quantification over a proper-class coefficient
space or use of a proper-class compactness theorem is required.
\end{corollary}
The assertion concerns natural-valuation precision of finite algebraic
data, not integration, a global surreal exponential, or ordinary paths in
the fine topology.

\section{Finite moment algebra and the cofactor coordinates}
\label{prony:sec:algebra}

Define
\begin{equation}\label{prony:eq:cofactor}
 P(X)=\prod_{i=1}^n(X-a_i),\quad
 Q_i(X)=\frac{P(X)}{X-a_i},\quad
 p_i=Q_i(a_i)=P'(a_i),\quad c_i=w_i p_i^2.
\end{equation}
Then $v(p_i)=d_i$ and $v(c_i)=E_i$. Both $P$ and every $Q_i$ are monic
integral polynomials. The $Q_i$ form a basis of $K[X]_{<n}$, since
evaluation at $a_j$ gives $Q_i(a_j)=\delta_{ij}p_i$.

Let $L:K[X]_{<2n}\to K$ be the linear functional with $L(X^k)=m_k$.
The moment representation is equivalent to
\[
 L(f)=\sum_i w_i f(a_i),\qquad \deg f<2n.
\]
In particular,
\begin{equation}\label{prony:eq:orthog}
 L(Pq)=0\quad(\deg q<n),\qquad
 L(Q_iQ_j)=\delta_{ij}c_i.
\end{equation}

\begin{proposition}[Classical Prony uniqueness]\label{prony:prop:prony}
Let $H=(m_{r+s})_{0\le r,s<n}$. Then
\begin{equation}\label{prony:eq:hankel}
 H=V\diag(w_i)V^{\mathsf T},\qquad
 \det H=\left(\prod_iw_i\right)\prod_{i<j}(a_j-a_i)^2\ne0,
\end{equation}
where $V_{ri}=a_i^r$. A regular $n$-node moment vector determines $P$
uniquely by its $n$ orthogonality equations, then the nodes up to permutation,
and then the weights. Any $n$-node representation of these moments is
automatically regular. There is no representation with fewer than $n$ nodes,
even if only $m_0,\ldots,m_{2n-2}$ are prescribed.
\end{proposition}
\begin{proof}
The displayed matrix factorization follows entry by entry. The determinant
uses the Vandermonde determinant. Writing
$P=X^n+\sum_{j<n}P_jX^j$, orthogonality against $1,X,\ldots,X^{n-1}$
is the linear system
\[
 H(P_0,\ldots,P_{n-1})^{\mathsf T}
 =-(m_n,\ldots,m_{2n-1})^{\mathsf T}.
\]
For any $n$-node representation, the same factorization of the invertible
$H$ forces its Vandermonde matrix and weight diagonal to be invertible.
Thus its nodes are distinct and its weights nonzero. It must give the
same monic annihilator and hence the same nodes. The first $n$ moments then determine
the weights through $V$. A configuration with fewer than $n$ nodes would
factor $H$ through a vector space of smaller dimension. This contradiction
uses only the moments through $m_{2n-2}$, which already determine $H$.
\end{proof}
Since $P$ is monic, \eqref{prony:eq:hankel} reads
$\det H=\bigl(\prod_iw_i\bigr)\Disc P$. With unit weights $H$ is the Hankel
matrix of the Newton power sums of $P$; this classical identity is a
common input shared with the Hankel section of the collection's spectral-theory
report (Section~\ref{prony:subsec:collection}).

\subsection{The structured correction system}
Let $\Delta L(X^k)=\epsilon_k$, and $\Lh=L+\Delta L$. Seek a monic
annihilator in the cofactor coordinates
\begin{equation}\label{prony:eq:Ph}
 \Ph=P+\sum_j b_jQ_j.
\end{equation}
Testing $\Lh(\Ph q)=0$ on $q=Q_i$ yields
\[
 c_i b_i+\sum_j\Delta L(Q_iQ_j)b_j=-\Delta L(PQ_i).
\]
Set
\begin{equation}\label{prony:eq:Bsystem}
 u_i=c_i b_i,\qquad
 B_{ij}=\frac{\Delta L(Q_iQ_j)}{c_j},\qquad
 g_i=\Delta L(PQ_i).
 \qquad (\Id+B)u=-g.
\end{equation}
This identity is exact, not a first-order approximation. In the $Q_i$
basis the perturbed Hankel bilinear form is
$(\Id+B)\diag(c_i)$. Thus invertibility of $\Id+B$ also proves invertibility
of the perturbed Hankel matrix.

\begin{lemma}[Uniform estimate in cofactor coordinates]\label{prony:lem:cofactorbound}
If $v(\epsilon_k)\ge\kappa$ and $\kappa>E_j$ for every $j$, system
\eqref{prony:eq:Bsystem} has a unique solution and
\begin{equation}\label{prony:eq:bprecision}
 v(u_i)\ge\kappa,\qquad v(b_i)\ge\kappa-E_i.
\end{equation}
\end{lemma}
\begin{proof}
Integrality of $Q_iQ_j$ and $PQ_i$ gives
$v(B_{ij})\ge\kappa-E_j>0$ and $v(g_i)\ge\kappa$.
The matrix $\Id+B$ has residue identity, so its inverse has entries in
$\Ocal$, by the adjugate formula. Multiplication by that inverse proves the
bounds. Alternatively its Neumann series is strongly summable by
Lemma~\ref{prony:lem:support}; either proof is valid at arbitrary rank.
\end{proof}

\section{A root-localization lemma adapted to the moment problem}
\label{prony:sec:roots}

\begin{lemma}[Cofactor localization]\label{prony:lem:localroots}
Suppose $v(b_i)>\delta_i$ for every $i$, and define $\Ph$ by
\eqref{prony:eq:Ph}. Then $\Ph$ has one and only one root $\ah_i$ in
$B_i=\{x:v(x-a_i)>\delta_i\}$, for every $i$. These are all its roots, and they
are simple. Writing $z_i=\ah_i-a_i$, one has
\begin{equation}\label{prony:eq:rootexact}
 z_i=-b_i\left(1+\sum_{j\ne i}
          \frac{b_j}{a_i-a_j+z_i}\right)^{-1}.
\end{equation}
Thus $z_i=0$ precisely when $b_i=0$, and otherwise
$v(z_i)=v(b_i)$ and $z_i/b_i\in-1+\mfr$.
\end{lemma}
\begin{proof}
For $X=a_i+z$ in $B_i$, none of the factors of $Q_i(X)$ vanishes. Dividing
$\Ph(X)$ by $Q_i(X)$ gives
\begin{equation}\label{prony:eq:localeq}
 z+b_i+\sum_{j\ne i}b_j\frac{z}{a_i-a_j+z}=0.
\end{equation}
Put $\tau=t^{\delta_i}$, $z=\tau y$, $D_{ij}=a_i-a_j$. The quantities
$\tau/D_{ij}$ are integral and $b_j/D_{ij}$ are infinitesimal, since
$\delta_j\ge d_{ij}$. Multiplying the scaled equation by
$\prod_{j\ne i}(1+(\tau/D_{ij})y)$ gives an integral polynomial whose
residue has value zero and derivative one at $y=0$. By
Lemma~\ref{prony:lem:hensel} it has exactly one root in $\mfr$. Its denominators
are units there, so it solves \eqref{prony:eq:localeq}.

The balls $B_i$ are disjoint: if a point belonged to $B_i$ and $B_j$, the
ultrametric inequality would give
$d_{ij}>\min(\delta_i,\delta_j)\ge d_{ij}$. Consequently the $n$ constructed roots
are distinct. A monic polynomial of degree $n$ with $n$ distinct roots
has precisely those roots, all simple. Every summand in the parentheses
in \eqref{prony:eq:rootexact} has positive valuation. The remaining assertions
follow from \eqref{prony:eq:localeq}.
\end{proof}

For later use, the lemma also gives
\begin{equation}\label{prony:eq:separation}
 \frac{\ah_i-\ah_j}{a_i-a_j}\in1+\mfr\qquad(i\ne j).
\end{equation}
Indeed both node displacements have valuation strictly greater than
$d_{ij}$. In particular, both the valuation and the complex leading
coefficient of each difference are unchanged.

\section{Reconstruction and the weight estimate}
\label{prony:sec:mainproof}

The node estimate alone is not the entire inverse problem. Estimating the
last Vandermonde solve as an unrelated operation generally loses the
correlations between the polynomial and moment errors. A cross-numerator
identity retains precisely those correlations.

\subsection{Recovering the numerator and all measured moments}
For a formal Laurent series in the independent variable $X^{-1}$, write
$[\cdot]_+$ for its polynomial part. This operation involves only finitely
many coefficients in the following formulas. Set
\begin{equation}\label{prony:eq:numerators}
 A=\sum_iw_iQ_i,
 \qquad
 \Ah=\left[\Ph(X)\sum_{k=0}^{2n-1}\mh_kX^{-k-1}\right]_+.
\end{equation}
Both numerators have degree less than $n$.

\begin{lemma}[Finite Pad\'e identity]\label{prony:lem:pade}
If $\Lh(\Ph q)=0$ for every $\deg q<n$, then
\begin{equation}\label{prony:eq:pade}
 \frac{\Ah(X)}{\Ph(X)}
 =\sum_{k=0}^{2n-1}\mh_kX^{-k-1}+O_X(X^{-2n-1}).
\end{equation}
When $\Ph$ has distinct roots $\ah_i$, its residues
$\wh_i=\Ah(\ah_i)/\Ph'(\ah_i)$ realize all $2n$ moments.
Here $O_X$ refers only to formal degree in $X^{-1}$, not to a valuation limit.
\end{lemma}
\begin{proof}
In the product used in \eqref{prony:eq:numerators}, the coefficient of $X^{-r-1}$,
for $0\le r<n$, is
$\sum_{j=0}^n\Ph_j\mh_{j+r}=\Lh(X^r\Ph)$, hence zero.
After subtracting the polynomial part, the product therefore begins at
$X^{-n-1}$. Division by the monic degree-$n$ polynomial produces
\eqref{prony:eq:pade}. Partial fractions give
$\Ah/\Ph=\sum_i\wh_i/(X-\ah_i)$, and expanding each fraction at infinity
identifies its first $2n$ coefficients.
\end{proof}

\begin{lemma}[Cross-numerator control]\label{prony:lem:cross}
Suppose $P,\Ph$ are monic integral degree-$n$ polynomials and their associated
rational functions agree with the original and perturbed moment vectors
through $X^{-2n}$. Then
\begin{equation}\label{prony:eq:cross}
 D=\Ah P-A\Ph,\qquad \deg D<2n,
\end{equation}
has every coefficient of valuation at least $\kappa$ whenever
$v(\epsilon_k)\ge\kappa$. In the cofactor coordinates,
\begin{equation}\label{prony:eq:Dvalues}
 D(a_i)=-w_i p_i^2b_i=-u_i.
\end{equation}
\end{lemma}
\begin{proof}
Multiply the Laurent expansion of $\Ah/\Ph-A/P$ by $P\Ph$ and take its
polynomial part. The tail $O_X(X^{-2n-1})$ contributes no polynomial term.
Every coefficient of $D$ is consequently a finite sum of moment errors
multiplied by integral coefficients of $P\Ph$. This proves the bound.
At $a_i$, use $A(a_i)=w_i p_i$ and $\Ph(a_i)=b_i p_i$.
\end{proof}

\subsection{A divided-difference estimate for the weights}
\begin{lemma}[Weight errors without a second conditioning loss]
\label{prony:lem:weight}
Under the hypotheses of Lemma~\ref{prony:lem:cross}, suppose the roots are labelled
by $v(\ah_i-a_i)>\delta_i$. Then
\[
 v(\wh_i-w_i)\ge\kappa-2d_i-\delta_i.
\]
\end{lemma}
\begin{proof}
Fix $i$, write $a=a_i$, $b=\ah_i$, and put
\[
 q(X)=Q_i(X)\frac{\Ph(X)}{X-\ah_i},\qquad F(X)=\frac{D(X)}{q(X)}.
\]
All factors of $q$ are nonzero at $a$ and $b$. Preservation of pairwise
separation gives
\begin{equation}\label{prony:eq:qval}
 v(q(a))=v(q(b))=2d_i.
\end{equation}
If $a\ne b$, the residues at $a$ and $b$ in
\[
 \frac{\Ah}{\Ph}-\frac A P=\frac{F(X)}{(X-a)(X-b)}
\]
are $-w_i$ and $\wh_i$. Thus
\begin{equation}\label{prony:eq:weightdd}
 \wh_i-w_i=\frac{F(b)-F(a)}{b-a}.
\end{equation}
If $a=b$, the left side has at most a simple pole at $a$, so $D(a)=0$,
and the residue gives $\wh_i-w_i=F'(a)$. This is the coincident-argument
version of the same formula.

For distinct arguments, denote a divided difference by $[a,b]f$.
Because $a,b\in\Ocal$ and every coefficient of $D$ has valuation at least
$\kappa$, expansion of $[a,b]X^k$ as a finite geometric sum gives
$v([a,b]D)\ge\kappa$. The same holds for $D'(a)$.
Each of the $2n-2$ linear factors of $q$ has value of valuation $d_{ij}$
at both arguments, with each $j\ne i$ occurring twice. Telescoping the
product difference, or differentiating in the coincident case, therefore
shows
\[
 v([a,b]q)\ge 2d_i-\delta_i.
\]
For $n=1$, $q=1$ and this divided difference is zero, which also satisfies
the bound. Finally the quotient identity
\[
 [a,b]F=\frac{[a,b]D}{q(b)}
       -\frac{D(a)[a,b]q}{q(a)q(b)}
\]
shows that its valuation is at least
$\min(\kappa-2d_i,\kappa-2d_i-\delta_i)=\kappa-2d_i-\delta_i$.
The derivative quotient formula proves the same estimate when $a=b$.
\end{proof}

\subsection{Proof of the reconstruction assertions}
\begin{proof}[Proof of Theorem~\ref{prony:thm:main}, except sharpness]
Since $\delta_i\ge0$ and $\kappa>\Theta$, Lemma~\ref{prony:lem:cofactorbound}
constructs the unique cofactor correction with
$v(b_i)\ge\kappa-E_i>\delta_i$. Lemma~\ref{prony:lem:localroots} constructs the $n$
simple roots, establishes the labelling, and proves \eqref{prony:eq:nodebound}.
They are integral, so their product $\Ph$ is integral.

Use Lemma~\ref{prony:lem:pade} to recover the weights and all measured moments.
Lemmas~\ref{prony:lem:cross} and \ref{prony:lem:weight} give
\eqref{prony:eq:weightbound}. Since this bound is strictly greater than
$v(w_i)=\varpi_i$, the new weights are nonzero and
$\wh_i/w_i\in1+\mfr$. Thus the constructed realization is regular.
The perturbed Hankel matrix is invertible by \eqref{prony:eq:Bsystem}; alternatively,
its factorization through these distinct nodes and nonzero weights proves
the same fact. Proposition~\ref{prony:prop:prony} gives uniqueness among all
regular $n$-node configurations. The balls are disjoint, so their labelling
is unique. Formula~\eqref{prony:eq:separation} supplies the final leading-term
assertion for the node differences.
\end{proof}

\begin{corollary}[Real and positive data]\label{prony:cor:positive}
For real Hahn or actual surreal inputs, the recovered nodes and weights
in Theorem~\ref{prony:thm:main} are real. Every initially positive weight remains
positive. The ordering of the nodes is preserved.
\end{corollary}
\begin{proof}
The construction takes place over the real coefficient field. Alternatively,
complex conjugation fixes the data and the labelled balls, so uniqueness
fixes every output. The ratios of corresponding weights and pairwise node
differences lie in $1+\mfr$, which consists of positive real Hahn elements.
\end{proof}

\begin{remark}[What the theorem does not require]
Weights need not be finite, positive, or have distinct leading coefficients.
Their valuations may be negative. Nodes may have identical standard parts,
and their mutual distances may have different Archimedean orders. No
assumption on the discriminant being an ordinary nonzero number is made.
The main proof uses only finite valued-field algebra and residue-simple
lifting; the Hahn setting additionally supplies explicit support control.
\end{remark}

\section{An exact worst direction and sharpness at every configuration}
\label{prony:sec:sharp}

The same perturbation direction works for every configuration: change only
the highest measured moment.
It is the familiar last-coordinate Prony-curve direction \cite{AGY}; the
point here is its exact threshold consequence.

\begin{proposition}[The last-moment formula]\label{prony:prop:last}
Let $\epsilon_{2n-1}=e$ and all other $\epsilon_k=0$. The original Hankel
matrix is unchanged. The unique monic annihilator determined by the new
moments is
\begin{equation}\label{prony:eq:Pe}
 P_e(X)=P(X)-e\sum_{i=1}^n\frac{Q_i(X)}{w_iP'(a_i)^2}.
\end{equation}
This formula is exact for every $e\in K$, regardless of its valuation.
\end{proposition}
\begin{proof}
Every $Q_iQ_j$ has degree at most $2n-2$, so $B=0$ in
\eqref{prony:eq:Bsystem}. Each $PQ_i$ is monic of degree $2n-1$, so $g_i=e$.
Thus $u_i=-e$ and $b_i=-e/c_i$. The unchanged Hankel matrix is invertible,
which proves uniqueness without any smallness hypothesis.
\end{proof}

\begin{proof}[Proof of the sharpness assertions in Theorem~\ref{prony:thm:main}]
If $v(e)=\kappa>\Theta$, then $v(b_i)=\kappa-E_i>\delta_i$.
Lemma~\ref{prony:lem:localroots} yields
\begin{equation}\label{prony:eq:leadingnode}
 v(\ah_i-a_i)=\kappa-E_i,\qquad
 \frac{\ah_i-a_i}{e/(w_iP'(a_i)^2)}\in1+\mfr.
\end{equation}
This proves simultaneous attainment of all the node bounds.

Now take $v(e)=\Theta$ and choose $i$ with $E_i+\delta_i=\Theta$.
Any regular $n$-node realization must have annihilator $P_e$, by
Proposition~\ref{prony:prop:prony}. Formula~\eqref{prony:eq:Pe} gives
\[
 v(P_e(a_i))=v(b_i p_i)=\Theta-E_i+d_i=\delta_i+d_i.
\]
If all its nodes admitted the strict-ball labelling, factorization at $a_i$
would instead give
\[
 v(P_e(a_i))
 =v(a_i-\ah_i)+\sum_{j\ne i}v(a_i-\ah_j)>\delta_i+d_i.
\]
The other factors have precisely their original separation valuations by
the strict-ball assumption. This is a contradiction.
\end{proof}

Because this boundary perturbation is allowed in every error ball with
precision $\kappa\le\Theta$, the universal strict-ball reconstruction
guarantee holds \emph{if and only if} $\kappa>\Theta$.

\subsection{Cluster interpretation and comparison with a determinant}
When every weight has valuation zero and all $n$ nodes occupy one equally
spaced valuation cluster, $d_{ij}=s>0$, one has
\[
 d_i=(n-1)s,\qquad \Theta=(2n-1)s.
\]
The pointwise geometry matters when this model is not uniform. A node deep
inside a nested cluster has one separation contribution for each other
node, not merely $n-1$ copies of the nearest separation.
The determinant records only a sum:
\begin{equation}\label{prony:eq:detval}
 v(\det H)=\sum_i\varpi_i+2\sum_{i<j}d_{ij}
           =\sum_i\varpi_i+\sum_i d_i.
\end{equation}
This single number does not specify the vector $(E_i+\delta_i)_i$ and does not
replace the threshold or the individual error bounds.

\subsection{Comparison with the coefficient threshold}
\label{prony:subsec:coefficient}
This subsection was added when the manuscript joined the collection. It
compares Theorem~\ref{prony:thm:main} with the coefficient-perturbation
theorem \lbl{polynomial:thm:stability} of the polynomial-algebra report
\cite{RepoPoly}. It proves no new theorem; every step is an elementary
consequence of results stated above or in that report.

\begin{remark}[Moment input versus coefficient input]
\label{prony:rem:coefficient}
For $P=\prod_i(X-a_i)$, the root data of \lbl{polynomial:thm:stability}
are the $d_{ij}$, $d_i$ and $\delta_i$ of \eqref{prony:eq:geometry}, and its
threshold is $T(P)=\max_i(d_i+\delta_i)$. By \eqref{prony:eq:geometry},
\[
 \Theta=\max_i\bigl((\varpi_i+d_i)+(d_i+\delta_i)\bigr).
\]
The extra summand $\varpi_i+d_i=v(w_iP'(a_i))$ is the cost, at $a_i$, of
passing from moments to the annihilator. Indeed, the perturbed annihilator
\eqref{prony:eq:Ph} differs from $P$ by $\sum_jb_jQ_j$, a polynomial of
degree less than $n$ whose value at $a_i$ is $b_ip_i$; every polynomial of
degree less than $n$ has this form, with $b_i$ its value at $a_i$ divided by
$p_i$. Lemma~\ref{prony:lem:cofactorbound} gives
$v(b_ip_i)\ge\kappa-\varpi_i-d_i$, and Lemma~\ref{prony:lem:localroots} needs
$v(b_ip_i)>d_i+\delta_i$, the $i$th summand of $T(P)$, at every node.
Replacing the coefficient precision $\varepsilon$ of
\lbl{polynomial:thm:stability} by the nodewise annihilator precision
$\kappa-\varpi_i-d_i$ turns its displacement bound $\varepsilon-d_i$ into the
node bound $\kappa-E_i$ of \eqref{prony:eq:nodebound}. If every $\varpi_i\ge0$,
then $\Theta\ge T(P)$. Weights of negative valuation can make $\Theta$
smaller: two nodes at separation valuation $s>0$ with both weights $t^{-2s}$
give $\Theta=s<2s=T(P)$.

Three further points make the comparison precise.
\begin{enumerate}[label=(\alph*)]
\item Lemma~\ref{prony:lem:localroots} is a nodewise form of the matching step
of \lbl{polynomial:thm:stability}. If every coefficient of $\Ph-P$ has
valuation greater than $T(P)$, then its values at the integral nodes satisfy
$v(b_ip_i)>d_i+\delta_i$, which is the hypothesis of the lemma. The converse
fails. For the nodes $0,t^s,1$ with $s>0$ one has $T(P)=2s$ and
$\delta_3=0$, so $\Ph=P+t^sQ_3$ satisfies the hypothesis of the lemma, while
$\Ph-P=t^sX^2-t^{2s}X$ has a coefficient of valuation $s<T(P)$. The lemma
yields the matching, the uniqueness of the labelling, the preserved
separations \eqref{prony:eq:separation} and the exact displacement valuation
$v(\ah_i-a_i)=v(b_i)$. It does not restate the second-order Newton estimate
or the discriminant statement of \lbl{polynomial:thm:stability}.
\item Theorem~\ref{prony:thm:main} does not follow merely by composing
Lemma~\ref{prony:lem:cofactorbound} with \lbl{polynomial:thm:stability},
which needs every coefficient of $\Ph-P$ to have valuation above $T(P)$. For
the two nodes $0,h=t^s$ with unit weights, the last-moment perturbation by
$e$ with $v(e)=\kappa$, formula \eqref{prony:eq:twoP}, gives $\Ph-P=-(2e/h^2)X+e/h$, with a coefficient of
valuation exactly $\kappa-2s$, while $T(P)=2s$. For $3s=\Theta<\kappa\le4s$
this perturbation is outside the hypothesis of
\lbl{polynomial:thm:stability}, whereas Theorem~\ref{prony:thm:main}
applies.
\item The sharpness statements differ in scope.
\lbl{polynomial:ex:sharpdisc} shows, for one family, that the strict
inequality in \lbl{polynomial:thm:stability} cannot be relaxed. Here the
last-moment direction of Proposition~\ref{prony:prop:last} defeats the
strict-ball guarantee at valuation $\Theta$ for every configuration. No
coefficient-side analogue of this universal statement is claimed.
\end{enumerate}
Finally, $\det H=\bigl(\prod_iw_i\bigr)\Disc P$, so \eqref{prony:eq:detval}
reads $v(\det H)=\sum_i\varpi_i+D$, where $D=v(\Disc P)=\sum_id_i$ is the
discriminant valuation of \lbl{polynomial:cor:discprecision}.
\end{remark}

\subsection{Infinite scales and affine normalization}
The integrality assumption on nodes is a normalization, not a restriction
on finite configurations in $\SC$. Choose $c,\rho\in K$, $\rho\ne0$, with
$x_i=(a_i-c)/\rho\in\Ocal$. The normalized moments are
\begin{equation}\label{prony:eq:affine}
 m^{\mathrm{norm}}_k
 =\rho^{-k}\sum_{\ell=0}^k\binom{k}{\ell}(-c)^{k-\ell}m_\ell.
\end{equation}
The same linear transformation must be applied to the errors. For error
bounds $v(\epsilon_\ell)\ge\kappa_\ell$, a valid normalized bound is
\[
 v(\epsilon^{\mathrm{norm}}_k)
 \ge\min_{0\le\ell\le k}
 \left(\kappa_\ell+(k-\ell)v(c)-k v(\rho)\right),
\]
with zero coefficients omitted; for $c=0$ only $\ell=k$ contributes.
Apply the uniform theorem to a common lower bound, or the nonuniform
certificate below to the individual normalized bounds. Translating back
adds $v(\rho)$ to each node-error valuation and does not change the weights.
One must not rescale the nodes while leaving the moment-error model
unchanged.

\section{Exact local precision and cancellation in the Hermite inverse}
\label{prony:sec:lattice}

\subsection{The differential and its classical dual basis}
Order the coordinates of the parameter space as
$(w_1,\ldots,w_n,a_1,\ldots,a_n)$, and write
$\mom:K^{2n}\to K^{2n}$ for the moment map \eqref{prony:eq:moments}.
At a regular configuration its differential is
\begin{equation}\label{prony:eq:diff}
 \dd m_k=\sum_i\left(a_i^k\dd w_i+w_i k a_i^{k-1}\dd a_i\right).
\end{equation}
The derivative term at $k=0$ is zero, including when $a_i=0$.
Define
\begin{equation}\label{prony:eq:hermite}
\begin{aligned}
 \ell_i(X)&=Q_i(X)/p_i,&
 \chi_i&=\ell_i'(a_i)=\sum_{j\ne i}\frac1{a_i-a_j},\\
 H_i(X)&=(1-2\chi_i(X-a_i))\ell_i(X)^2,&
 G_i(X)&=(X-a_i)\ell_i(X)^2.
\end{aligned}
\end{equation}
These are the usual degree-$<2n$ Hermite interpolation polynomials.
They satisfy
\[
 H_i(a_j)=\delta_{ij},\quad H_i'(a_j)=0,\qquad
 G_i(a_j)=0,\quad G_i'(a_j)=\delta_{ij}.
\]
Consequently the inverse differential is
\begin{equation}\label{prony:eq:invdiff}
 \dd w_i=(\dd L)(H_i),\qquad
 \dd a_i=(\dd L)(G_i)/w_i.
\end{equation}
This proves that the Jacobian $J$ is invertible and identifies each row of
$J^{-1}$. These interpolation identities are classical; the precision
calculation below specifies exactly what they imply over Hahn coefficients.

For a nonzero polynomial $f$, let
$\vg(f)=\min_k v([X^k]f)$. This Gauss valuation is multiplicative: normalize
each factor by a monomial to make its minimum zero and reduce modulo
$\mfr$; the product of two nonzero residue polynomials is nonzero.
Define the cancellation depth
\begin{equation}\label{prony:eq:cancellation}
 q_i=\begin{cases}\max(0,-v(\chi_i)),&\chi_i\ne0,\\0,&\chi_i=0.\end{cases}
 \qquad 0\le q_i\le \delta_i.
\end{equation}
The inequality follows from the ultrametric inequality for the sum defining
$\chi_i$.

\begin{proposition}[Exact row losses]\label{prony:prop:rows}
For the integral nodes fixed above,
\begin{equation}\label{prony:eq:gaussrows}
 \vg(G_i)=-2d_i,\qquad \vg(H_i)=-2d_i-q_i.
\end{equation}
Thus the exact coordinate projections of the tangent precision lattice
$J^{-1}t^\kappa\Ocal^{2n}$ are
\begin{equation}\label{prony:eq:marginals}
 \dd a_i\in t^{\kappa-E_i}\Ocal,
 \qquad
 \dd w_i\in t^{\kappa-2d_i-q_i}\Ocal.
\end{equation}
Each displayed projection is the whole indicated fractional ideal, not
merely a containing ball.
\end{proposition}
\begin{proof}
The monic integral polynomials $Q_i$ and $X-a_i$ have Gauss valuation zero,
so $\vg(\ell_i^2)=-2d_i$ and the first formula follows.
If $\chi_i=0$, the remaining linear factor in $H_i$ is one. If
$v(\chi_i)<0$, its coefficient of $X$ has valuation $v(\chi_i)$ and its constant
coefficient has at least that valuation. If $v(\chi_i)\ge0$, the factor is
integral and evaluates to one at the integral point $a_i$, so its Gauss
valuation is zero. Here $v(2)=0$ since the coefficient field is $\R$ or
$\C$. Multiplicativity gives the second formula.

The image of an $\Ocal$-linear row is the fractional ideal generated by its
finitely many coefficients. Such an ideal is generated by a coefficient
of minimum valuation. Equations~\eqref{prony:eq:invdiff} and
\eqref{prony:eq:gaussrows} now prove \eqref{prony:eq:marginals}, including surjectivity
onto each ideal separately.
\end{proof}

\subsection{A support-based exact inverse-lattice theorem}
The following proposition is an arbitrary-rank polynomial implementation
of the differential precision principle in \cite{CRV}. It is included with
proof so that no Banach-space convergence hypothesis is silently imported.

\begin{proposition}[Finite coefficient criterion]\label{prony:prop:lattice}
Let $F:K^N\to K^N$ be polynomial, let $x\in K^N$, and suppose
$J=\dd F_x$ is invertible. Expand, in independent variables $Z$,
\begin{equation}\label{prony:eq:Cexpand}
 F(x+J^{-1}Z)-F(x)
 =Z+\left(\sum_{|\nu|\ge2} C_{j,\nu} Z^\nu\right)_{j=1}^N.
\end{equation}
The sums are finite. If $\kappa\in\Gamma$ satisfies
\begin{equation}\label{prony:eq:latticecriterion}
 v(C_{j,\nu})+(|\nu|-1)\kappa>0
 \quad\hbox{for every nonzero }C_{j,\nu},
\end{equation}
then $F$ is a bijection from $x+J^{-1}t^\kappa\Ocal^N$ onto
$F(x)+t^\kappa\Ocal^N$. In particular,
\begin{equation}\label{prony:eq:latticeimage}
 F(x+J^{-1}t^\kappa\Ocal^N)=F(x)+t^\kappa\Ocal^N.
\end{equation}
\end{proposition}
\begin{proof}
Scale $Z=t^\kappa Y$ and divide the output increment by $t^\kappa$.
The normalized map is
$T(Y)=Y+R(Y)$, where every coefficient of $R$ belongs to $\mfr$ and every
monomial has degree at least two.
For $Y,Y'\in\Ocal^N$, polynomial divided differences give
\[
 R(Y)-R(Y')=A(Y,Y')(Y-Y')
\]
with every entry of $A(Y,Y')$ in $\mfr$. If $Y\ne Y'$, the minimum
valuation of this correction is strictly greater than
$v_{\min}(Y-Y')$. It follows that
\[
 v_{\min}(T(Y)-T(Y'))=v_{\min}(Y-Y'),
\]
which proves injectivity.

For a prescribed $b\in\Ocal^N$, the polynomial system $T(Y)=b$ reduces
to $Y=\overline b$ modulo $\mfr$, with residue Jacobian identity.
Lemma~\ref{prony:lem:hensel} gives a solution in
$\overline b+\mfr^N\subset\Ocal^N$. This proves surjectivity.
The lemma's proof evaluates a formal series in finitely many positive-support
coefficients and in $b-\overline b$; no assertion about an ordinary Cauchy
sequence is involved. Undo the two linear changes of variables.
\end{proof}

\begin{corollary}[Exact nonlinear uncertainty set]\label{prony:cor:nonlinear}
For the Prony map, assume $\kappa>\Theta$ and also
\eqref{prony:eq:latticecriterion} for its finite expansion \eqref{prony:eq:Cexpand}.
Then the set of all parameter increments with strict-ball labelling realizing moment errors
in $t^\kappa\Ocal^{2n}$ is exactly
\begin{equation}\label{prony:eq:exactinverse}
 J^{-1}t^\kappa\Ocal^{2n}.
\end{equation}
In particular the exact coordinate ranges are
\[
 \ah_i-a_i\in t^{\kappa-E_i}\Ocal,
 \qquad
 \wh_i-w_i\in t^{\kappa-2d_i-q_i}\Ocal.
\]
\end{corollary}
\begin{proof}
By Proposition~\ref{prony:prop:rows}, every point of the proposed parameter
lattice satisfies strict node matching and preserves nonzero weights:
$q_i\le \delta_i$ and $\kappa>\Theta$ give the required inequalities.
Proposition~\ref{prony:prop:lattice} proves exact surjectivity onto the error
ball, and Theorem~\ref{prony:thm:main} gives uniqueness of its labelled
realization. The coordinate projections were computed exactly in
Proposition~\ref{prony:prop:rows}.
\end{proof}

The lattice in \eqref{prony:eq:exactinverse} is generally not the Cartesian
product of its coordinate projections: errors in different recovered
parameters are correlated. Since $\Bbbk$ is infinite, all finite row minima
can nevertheless be attained simultaneously. Normalize each nonzero row
of $J^{-1}$ by its least valuation and choose an ordinary vector avoiding
the finitely many resulting residue hyperplanes. Its lattice image has
the minimum valuation in every parameter coordinate, and its actual moment
increment belongs to the prescribed error ball.

There always exists a sufficiently large $\kappa$ satisfying the finite
criterion: choose it positive and larger than every $-v(C_{j,\nu})$ and
than $\Theta$. In a divisible group one may use the sharper sufficient
inequalities
$\kappa>-v(C_{j,\nu})/(|\nu|-1)$.
This auxiliary criterion is \emph{not} claimed to be the optimal onset of
exact lattice behavior. It does not weaken the sharp threshold theorem,
which requires only $\kappa>\Theta$ and was proved independently.

\section{Nonuniform moment precision and a directed-path certificate}
\label{prony:sec:graph}

Uniform measurement precision is convenient but sometimes wasteful.
Keep the original nontrivial ordered abelian group $\Gamma$; divisibility
is not required. Suppose
\[
 v(\epsilon_k)\ge\kappa_k\in\Gamma,\qquad 0\le k<2n.
\]
For a nonzero polynomial of degree less than $2n$, define its noise gauge
\begin{equation}\label{prony:eq:noisegauge}
 \nu_{\boldsymbol\kappa}(f)
 =\min_{[X^k]f\ne0}\bigl(\kappa_k+v([X^k]f)\bigr).
\end{equation}
This always bounds $v(\Delta L(f))$ from below, although cancellations may
make the actual value higher. Put
\begin{equation}\label{prony:eq:graphedges}
 \lambda_i=\nu_{\boldsymbol\kappa}(PQ_i),\qquad
 \mu_{ij}=\nu_{\boldsymbol\kappa}(Q_iQ_j),\qquad
 g_{ij}=\mu_{ij}-E_j.
\end{equation}
Make a directed graph with vertices $1,\ldots,n$; the edge from $j$ to $i$
has weight $g_{ij}$. Loops are included. A path starting at $j$ has initial
cost $\lambda_j$, plus the sum of its edge weights.

\begin{theorem}[Nonuniform reconstruction certificate]\label{prony:thm:graph}
Suppose every directed simple cycle, including every loop, has strictly
positive total edge weight. For each $i$, let $\rho_i$ be the minimum cost
of a simple path ending at $i$, with length-zero paths allowed. If
\begin{equation}\label{prony:eq:graphcondition}
 \rho_i-E_i>\delta_i\quad\hbox{for every }i,
\end{equation}
then every allowed moment perturbation has a unique regular $n$-node
realization. It admits strict-ball labelling and satisfies
\begin{equation}\label{prony:eq:graphnode}
 v(\ah_i-a_i)\ge\rho_i-E_i.
\end{equation}
If $\kappa_* =\min_k\kappa_k$, a further, possibly weak, weight bound is
$v(\wh_i-w_i)\ge\kappa_*-2d_i-\delta_i$.
The certificate is sufficient; no necessity or universal sharpness is
claimed for these graph conditions.
\end{theorem}

\begin{proof}
Use the ordered divisible hull $\Gamma_{\Q}$ of $\Gamma$ temporarily.
It consists of fractions $\gamma/m$ with $\gamma\in\Gamma$ and positive
integers $m$; equality and order are defined by cross multiplication.
Since an ordered abelian group is torsion-free, $\Gamma$ embeds by
$\gamma\mapsto\gamma/1$. Keeping the same coefficients and supports gives
a valuation-preserving field embedding
\[
 K\subset K^{\mathrm{aux}}=\Bbbk((t^{\Gamma_{\Q}})).
\]
Choose $\eta\in\Gamma_{\Q}$ positive and smaller than every simple cycle's
mean edge weight, for instance half the minimum of the finitely many
positive means. Replace each edge weight by $g_{ij}-\eta$. Every directed
cycle now has positive weight, because a closed walk decomposes into simple
cycles. Add a source with a zero-weight edge to every vertex, and let
$\zeta_i$ be its shortest-path cost to $i$. Shortest costs are attained
by simple paths, since deleting cycles decreases cost. They satisfy
\[
 \zeta_i\le\zeta_j+g_{ij}-\eta,
 \qquad g_{ij}+\zeta_j-\zeta_i\ge\eta.
\]
Let $T=\diag(t^{\zeta_i})$. The matrix
$B'=T^{-1}BT$ has every entry of positive valuation, by
\eqref{prony:eq:Bsystem} and \eqref{prony:eq:graphedges}. Hence $\Id+B'$ and $\Id+B$ are invertible in the auxiliary field.
But $\det(\Id+B)\in K$ is now known to be nonzero, so the adjugate formula
places $(\Id+B)^{-1}$ in $K$ itself. Moreover the Neumann expansion is strongly summable,
by Lemma~\ref{prony:lem:support} applied to the finitely many entries of $B'$.
Multiplying by a fixed finite vector of Hahn series preserves this strong
summability. Since $B^\ell=T(B')^\ell T^{-1}$, the unscaled walk terms
have the same property: for each pair of endpoints their supports differ
by one fixed translation, and there are only finitely many endpoints.

In
\[
 u=-\sum_{\ell\ge0}(-B)^\ell g,
\]
every matrix-product term is indexed by a directed walk ending at its
output vertex. Its valuation is at least the initial $\lambda$ plus its
edge weights. Deleting repeated-vertex cycles only decreases this lower
bound. Thus every such term has valuation at least $\rho_i$, and so
$v(u_i)\ge\rho_i$. Strong summability legitimizes the expansion and
regrouping without requiring the multiples of $\eta$ to be cofinal.
Every term of the unscaled expansion has support in $\Gamma$, so its
sum also belongs to $K$: well-ordering and finite multiplicity in the
auxiliary ordered group restrict to the original subgroup. The finite
simple-path costs $\rho_i$ lie in $\Gamma$ too. It follows that the solution
in $K$ satisfies $v(b_i)\ge\rho_i-E_i>\delta_i$.

Construct the $n$ roots in the original field $K$ by Lemma~\ref{prony:lem:localroots} and the weights by
Lemma~\ref{prony:lem:pade}. These weights are nonzero even when the coarse weight
bound is inconclusive: the perturbed Hankel matrix is invertible because
its matrix in the $Q_i$ basis is $(\Id+B)\diag(c_i)$, whereas its
factorization through the constructed distinct nodes would be singular
if any weight vanished. Classical Prony uniqueness completes the
reconstruction assertion. The root valuation follows from the same lemma;
Lemma~\ref{prony:lem:weight} with $\kappa_*$ gives the last bound.
\end{proof}

Fractional potentials can be needed even when all edge weights are
integers. For two vertices with off-diagonal edge weights $-1$ and $2$
and loops of weight $1$, all simple cycles are positive. Integer
potentials cannot make both off-diagonal weights positive: they would
then each be at least $1$, whereas their sum is invariant and equals $1$.
The auxiliary divisible hull supplies the scaling without changing the
field of the input or reconstructed output.

\subsection{How to compute the certificate}
The certificate requires finite polynomial multiplication, finitely many
coefficient valuations, and ordered-group arithmetic. One may enumerate
simple cycles and paths for a small configuration. More efficiently,
shortest-path relaxation supplies the potentials and the costs once a
valid positive margin is chosen. The value group need not be represented
by real numbers: rational tuples with lexicographic comparison work directly.
The polynomial coefficients themselves must be known well enough to decide
zero coefficients and their valuations.

The gauge \eqref{prony:eq:noisegauge} discards some information: the quantities
$\Delta L(PQ_i)$ and $\Delta L(Q_iQ_j)$ are built from the \emph{same}
errors. Cancellations and cross-entry relations can make the graph
certificate fail while reconstruction is still well behaved. It would be
incorrect to advertise this sufficient method as the exact admissible
precision region. In the uniform case, taking
$\lambda_i\ge\kappa$ and $g_{ij}\ge\kappa-E_j>0$ immediately recovers the
bounds used in Theorem~\ref{prony:thm:main}.

\section{Worked examples}
\label{prony:sec:examples}

\subsection{One atom and the meaning of the empty maximum}
For $n=1$, $m_0=w$ and $m_1=wa$. The exact perturbed parameters are
\[
 \wh=w+\epsilon_0,\qquad
 \ah=\frac{wa+\epsilon_1}{w+\epsilon_0},\qquad
 \ah-a=\frac{\epsilon_1-a\epsilon_0}{w+\epsilon_0}.
\]
Our convention gives $d=\delta=0$, $E=\Theta=\varpi=v(w)$.
Errors of valuation at least $\kappa>\varpi$ preserve the weight's leading
term and put $\ah$ in $a+\mfr$, with node valuation bound
$\kappa-\varpi$ and weight bound $\kappa$. Taking
$\epsilon_0=0$, $\epsilon_1=e$ attains the node bound; at
$v(e)=\varpi$, the new node is not in $a+\mfr$.
This example checks the degree-one endpoint conventions directly.

\subsection{Two nodes: exact formulas and an actual boundary collision}
Let $a_1=0$, $a_2=h=t^s$, $s>0$, and $w_1=w_2=1$. Then
\[
 (m_0,m_1,m_2,m_3)=(2,h,h^2,h^3),\quad
 d_i=\delta_i=s,\quad E_i=2s,\quad\Theta=3s.
\]
Perturb only $m_3$ by $e$. Proposition~\ref{prony:prop:last} gives
\begin{equation}\label{prony:eq:twoP}
 P_e(X)=X^2-\left(h+\frac{2e}{h^2}\right)X+\frac eh,
 \qquad \Disc(P_e)=h^2+\frac{4e^2}{h^4}.
\end{equation}
For $e=h^4$, put $S=\sqrt{1+4h^2}$, the root in $1+\mfr$.
The strict-ball-labelled nodes and weights are
\begin{align*}
 \ah_1&=(h+2h^2-hS)/2,&\wh_1&=1+2h/S,\\
 \ah_2&=(h+2h^2+hS)/2,&\wh_2&=1-2h/S.
\end{align*}
Their initial terms are
\begin{align*}
 \ah_1&=h^2-h^3+h^5-2h^7+O_h(h^9),\\
 \ah_2&=h+h^2+h^3-h^5+2h^7+O_h(h^9),\\
 \wh_1&=1+2h-4h^3+12h^5-40h^7+O_h(h^9),\\
 \wh_2&=1-2h+4h^3-12h^5+40h^7+O_h(h^9).
\end{align*}
Here $O_h(h^9)$ denotes formal power-series order, hence Hahn valuation
at least $9s$. Both node displacements have valuation $2s=4s-E_i$;
both weight errors have valuation $s=4s-2d_i-\delta_i$.

At the boundary take $e=c h^3$ with ordinary $c\ne0$. Formula
\eqref{prony:eq:twoP} becomes
\[
 P_e(X)=X^2-h(1+2c)X+ch^2.
\]
For $c=i/2$, the discriminant is zero and the root is
$h(1+i)/2$. The unchanged Hankel determinant is $h^2\ne0$, so no regular
two-node realization exists. For nonzero real $c$, the roots instead are
\[
 \frac h2\left(1+2c\pm\sqrt{1+4c^2}\right),
\]
with positive weights
$1\pm2c/\sqrt{1+4c^2}$ assigned in the reverse sign order.
They give a regular positive real realization, but neither node retains
its strict nearest-neighbour label. Thus failure at the threshold need
not mean a collision or a loss of positivity.

\subsection{The same tree, different weight precision}
Let the three nodes be $-h,0,h$, with $h=t^s$, $s>0$, and all weights one.
Every separation valuation is $s$, so
\[
 d_i=2s,\quad \delta_i=s,\quad E_i=4s,\quad\Theta=5s.
\]
At the middle node,
\[
 \chi_2=\frac1h+\frac1{-h}=0,
\]
whereas at the two endpoints $\chi_1=-3/(2h)$ and $\chi_3=3/(2h)$.
Hence $q_2=0$ and $q_1=q_3=s$.
The exact weight row losses are consequently $4s$ at the middle and $5s$
at the endpoints. Under the explicit extra smallness criterion of
Corollary~\ref{prony:cor:nonlinear}, the middle weight's full nonlinear uncertainty
range is $t^{\kappa-4s}\Ocal$, a factor $t^s$ smaller than the endpoint
ranges $t^{\kappa-5s}\Ocal$.

For comparison take the nodes $0,h,3h$. They have exactly the same labelled
separation valuations. Now
\[
 \chi_1=-\frac4{3h},\qquad \chi_2=\frac1{2h},\qquad \chi_3=\frac5{6h},
\]
so all three $q_i$ equal $s$. Thus the full separation matrix determines
the uniform threshold and the node losses, but does \emph{not} determine
the exact weight precision. Residue-level cancellation contains additional
information. The extra nonlinear lattice criterion is essential to the
stated nonlinear interpretation; the tangent calculation needs no such
criterion.


\begin{proposition}[A nonlinear obstruction to the cancellation gain]
\label{prony:prop:nonlinearobstruction}
For the symmetric configuration $(-h,0,h)$ with unit weights, perturb only
$m_5$ by $e$, with $v(e)=\kappa>5s$. The recovered middle weight satisfies
\begin{equation}\label{prony:eq:nonlinearweight}
 v(\wh_2-1)=2\kappa-10s.
\end{equation}
Consequently, whenever $5s<\kappa<6s$, its nonlinear error violates the
improved tangent bound $\kappa-4s$. The exact Hermite lattice cannot describe
the whole inverse error set throughout the entire sharp-threshold region.
\end{proposition}
\begin{proof}
Normalize the nodes by $X=hY$ and put $\lambda=e/h^5\in\mfr$.
The original normalized moments are $(3,0,2,0,2,0)$ and the perturbed last
moment is $\lambda$. The exact monic annihilator and numerator are
\[
 P_\lambda(Y)=Y^3-\frac32\lambda Y^2-Y+\lambda,
 \qquad
 A_\lambda(Y)=3Y^2-\frac92\lambda Y-1.
\]
These follow from \eqref{prony:eq:Pe} and \eqref{prony:eq:numerators}, or directly
from the three orthogonality equations. The middle root $y\in\mfr$
satisfies
\[
 y\left(1-y^2+\frac32\lambda y\right)=\lambda,
\]
so $v(y)=v(\lambda)=\kappa-5s$. Its recovered weight obeys
\[
 \wh_2-1
 =\frac{A_\lambda(y)-P_\lambda'(y)}{P_\lambda'(y)}
 =\frac{-\tfrac32\lambda y}{3y^2-3\lambda y-1}.
\]
The denominator is a unit of residue $-1$, giving
\eqref{prony:eq:nonlinearweight}. For $\kappa<6s$,
$2\kappa-10s<\kappa-4s$, as asserted. Values in this interval exist in
the divisible workspaces used for actual surreal numbers.
\end{proof}

For the same symmetric example, the explicit choice $\kappa>9s$ is a
simple sufficient condition for the full nonlinear lattice theorem.
Indeed the inverse Hermite lattice gives weight increments of valuation
at least $\kappa-5s$ and node increments of valuation at least
$\kappa-4s$. In the $k$th moment, a nonlinear term with $j\ge2$ node
increments and no weight increment has valuation at least
$j\kappa+(k-5j)s$, with $k\ge j$. A term with one weight increment and
$j\ge1$ node increments has valuation at least
$(j+1)\kappa+(k-5j-5)s$, again with $k\ge j$.
For $\kappa>9s$ both bounds are strictly greater than $\kappa$.
The zeroth moment is linear. Applying these estimates coefficientwise in
the independent normalized moment variables verifies
\eqref{prony:eq:latticecriterion}. This is a sufficient hand-checkable bound,
not an optimal one; the obstruction just proved occurs below $6s$.

\subsection{A genuinely higher-rank configuration}
Take $\Gamma=\Q e_1\oplus_{\mathrm{lex}}\Q e_2$, with
$e_1>Ne_2>0$ for every positive integer $N$, and put
\[
 (a_1,a_2,a_3,a_4)=(0,t^{e_1},t^{e_2},1),\qquad
 (w_1,w_2,w_3,w_4)=(1,t^{e_2},1,t^{2e_2}).
\]
The only positive separation valuations are
$d_{12}=e_1$ and $d_{13}=d_{23}=e_2$.
The following table records the full threshold data.
\begin{center}
\begin{tabular}{c c c c c}
\toprule
$i$ & $\varpi_i$ & $d_i$ & $\delta_i$ & $E_i+\delta_i$\\
\midrule
1 & $0$ & $e_1+e_2$ & $e_1$ & $3e_1+2e_2$\\
2 & $e_2$ & $e_1+e_2$ & $e_1$ & $3e_1+3e_2$\\
3 & $0$ & $2e_2$ & $e_2$ & $5e_2$\\
4 & $2e_2$ & $0$ & $0$ & $2e_2$\\
\bottomrule
\end{tabular}
\end{center}
Thus $\Theta=3e_1+3e_2$. For $\kappa=3e_1+4e_2$, the guaranteed node
precisions are, respectively,
\[
 e_1+2e_2,\qquad e_1+e_2,\qquad 3e_1,\qquad 3e_1+2e_2.
\]
The weight-error bounds are
\[
 2e_2,\qquad 2e_2,\qquad 3e_1-e_2,\qquad 3e_1+4e_2.
\]
Each weight-error bound is strictly above the original weight's valuation.
Perturbing the seventh moment, $m_7$, by $t^\kappa$ attains all four node
bounds exactly. All these assertions are ordered-pair arithmetic; no
real-valued logarithm is involved.

In actual surreals one may choose $e_1=\omega$, $e_2=1$. The nodes then
are $0,\omega^{-\omega},\omega^{-1},1$, and the admissible test error is
$\omega^{-(3\omega+4)}$. There is no ordered-group embedding of this entire
$\Gamma$ into the additive reals: a positive image of $e_2$ cannot have
all its integer multiples below a fixed image of $e_1$. This does not deny
that a finite collection of these inequalities can sometimes be emulated
by a sufficiently large real scale ratio; it explains why one such ratio
cannot represent the whole group.

\subsection{Nonuniform errors for two nodes}
Return to $0,h$ with weights one. For $Q_1=X-h$, $Q_2=X$ and
$P=X^2-hX$, the noise gauges are
\begin{align*}
 \lambda_1&=\min(\kappa_3,\kappa_2+s,\kappa_1+2s),\\
 \lambda_2&=\min(\kappa_3,\kappa_2+s),\\
 \mu_{11}&=\min(\kappa_2,\kappa_1+s,\kappa_0+2s),\\
 \mu_{12}=\mu_{21}&=\min(\kappa_2,\kappa_1+s),\qquad
 \mu_{22}=\kappa_2.
\end{align*}
Suppose
\begin{equation}\label{prony:eq:hetero}
 \kappa_0>0,\qquad\kappa_1>s,\qquad
 \kappa_2>2s,\qquad\kappa_3>3s.
\end{equation}
Then all graph edges have positive weight, and all initial costs exceed
$3s$. Every $\rho_i$ therefore exceeds $3s$, so the certificate applies.
These errors need not have a common precision above $3s$.

There is also a simple direct interpretation: normalize $X=hY$ and divide
the $k$th moment by $h^k$. The normalized nodes are $0,1$, their threshold
is zero, and all normalized errors in \eqref{prony:eq:hetero} are infinitesimal.
The main theorem then also proves preservation of the leading weights.
This example illustrates the certificate; it is not claimed as a new
scaling principle.

\section{Constructive content, symbolic checks, and formalization}
\label{prony:sec:verification}

\subsection{A finite-algebra reconstruction procedure}
With exact input data and a certified precision hypothesis, the proof gives
a specific procedure. Form $P,Q_i,p_i,c_i$, compute the finite quantities
$\Delta L(Q_iQ_j)$ and $\Delta L(PQ_i)$, solve
$(\Id+B)u=-g$, and form $\Ph=P+\sum_i(u_i/c_i)Q_i$.
For each $i$, solve the residue-simple local equation
\eqref{prony:eq:localeq} by the formal-parameter construction. Finally compute
$\Ah$ by finite polynomial-part extraction and the weights by residues.
All node corrections and their labels are certified before the last step.

This is a mathematical construction, not an unconditional Turing algorithm
on arbitrary surreal inputs. An effective implementation needs names
supporting exact coefficient arithmetic, equality, leading valuations, and
the required support operations. Finite algebra does not by itself supply
those oracles. Over rational-function or bounded-denominator Puiseux inputs,
finite truncations can be computed using familiar exact symbolic methods;
those representations are narrower than the full Hahn field.

The formal lifting proofs are support-controlled in a precise relative
sense. After the finite normalizations, record every positive-support
parameter used in the inverse and residue-simple equations. The corresponding
formal solution supports lie in the monoid generated by those parameter
supports. Monomial rescalings add their explicitly known shifts. This is
\emph{not} an assertion that the output uses only the supports of the raw
moment errors: inverse separations, normalized units, and starting
coefficients contribute to the normalized parameters as well.

\subsection{What the accompanying program checks}
This directory contains \lbl{code/verify.py}, which uses exact SymPy
arithmetic and writes a human-readable verification log. Its tests cover:
\begin{enumerate}
\item the Hermite cardinal conditions and the matrix identity
      $J^{-1}J=\Id$ for several finite rational configurations;
\item the exact cofactor correction system, moment annihilator,
      cross-numerator identity, and recovered finite Laurent coefficients
      for explicit rational perturbations;
\item the two-node formula \eqref{prony:eq:twoP}, its discriminant and boundary
      collision, and the displayed power-series coefficients;
\item the three-node cancellation comparison, its nonlinear middle-weight
      obstruction, and the rank-two threshold arithmetic;
\item a directed graph with a negative edge but positive cycles, verifying
      shortest-path potentials and the path lower bounds.
\end{enumerate}
There is no floating-point evidence in these checks. They detect sign,
index, and algebraic mistakes in concrete instances. They do not prove
arbitrary-rank summability, the general quantified theorems, or any novelty
claim. Those rest on the mathematical arguments, with their stated
classical prerequisites.

The recorded run, \lbl{data/verification.txt}, passed $10{,}052$ exact
assertions under Python~3.13.5 and SymPy~1.14.0; the count includes finite
graph-walk enumeration and is not a count of independent results. The
program first deletes that file and then rewrites it, recording the Python
and SymPy versions on its second line, so it should be run on a copy of this
directory. When the report joined the collection, a run on such a copy under
Python~3.14.4 and SymPy~1.14.0 reproduced every line of the record except
that version line. The program keeps the manuscript's letters
(Section~\ref{prony:subsec:collection}), and
\lbl{data/requirements.txt} pins the tested SymPy version.

\subsection{Current Lean coverage and remaining development}
The ledger \lbl{docs/FORMALIZATION.md} records the finite algebra already
proved in \lbl{Surreal/Algebra/PronyHankel.lean}: classical Prony uniqueness,
the Hankel determinant, cofactor orthogonality, the exact last-moment
annihilator, and the algebraic local-root identities.
Valuation localization and the remaining precision estimates are pending.

The dependency order separates finite algebra from valuation lemmas.
The first layer completes the field-level algebra above with the finite
Pad\'e identity. The second layer adds a valuation, proves integral-polynomial
and divided-difference estimates, and packages residue-simple lifting as an
explicit assumption. It then proves the sharp uniform theorem without any
reference to a proper class.

The third layer instantiates the lifting assumption for Hahn fields via
strong formal evaluation, and transfers finite families to the actual
surreal carrier through the normal-form bridge. The Hermite inverse layer
can reuse the repository's finite interpolation machinery: the collection's
formalization ledger \lbl{docs/FORMALIZATION.md} maps the Hermite and
Lagrange interpolation clauses of \lbl{polynomial:thm:crt} to checked Lean
declarations over fields. The graph
certificate is finite ordered-group combinatorics plus matrix scaling.
The exact lattice theorem needs the multivariable residue-simple lifting
lemma, not an ordinary normed-space inverse-function theorem.

The build and axiom audit cover the imported declarations recorded in the
ledger; they do not establish the unformalized precision theorems.
Completing the mechanization requires matching each
hypothesis, particularly the exact model order, strict inequalities, and
universe sizes.

\section{Boundaries of the results and further questions}
\label{prony:sec:scope}

The uniform theorem resolves the precise finite reconstruction question
posed here, including optimality at each configuration. It is not presented
as a solution to a named previously published open problem. The primary
research claim is a proposed new exact refinement with complete proofs,
subject to independent verification and a broader priority search.

Three distinctions are particularly important. First, the threshold is
optimal for retaining all strict nearest-neighbour labels, not for the
bare existence of any atomic representation. Second, the weight bound
$\kappa-2d_i-\delta_i$ is uniformly valid above the sharp threshold, whereas the
improved cancellation bound with $q_i$ is proved for the tangent map and,
under a separate finite coefficient criterion, for the whole nonlinear
uncertainty set. Third, the directed-path theorem is a sufficient
nonuniform certificate, not the exact geometry of all allowed error ideals.

Several focused problems remain. An exact nonuniform precision region
would have to retain algebraic correlations among the entries of the
cofactor correction matrix, rather than just their valuations. A sharp
criterion for when the entire inverse image is its Hermite lattice may be
weaker than \eqref{prony:eq:latticecriterion}. Confluent configurations, where
several atoms merge and derivative evaluations replace simple evaluations,
require a different local model. Coincident nodes have a zero difference
and hence separation valuation $+\infty$; valuation zero instead describes
distinct residue classes and is already allowed here. Finally, infinite configurations need a
separate summability hypothesis on the moment data and are outside this
finite theorem.

These questions are stated as research directions, not asserted to be
known open problems in the published literature. The present finite
results already supply a useful answer: the information needed to resolve
surcomplex atoms is governed by a weighted, pointwise collision profile,
and the exact worst moment direction can be written down without solving
any nonlinear equation.

\appendix
\section{Dependency and novelty ledger}
\label{prony:app:ledger}

\begin{longtable}{>{\raggedright\arraybackslash}p{0.31\textwidth}>{\raggedright\arraybackslash}p{0.60\textwidth}}
\toprule
Statement or ingredient & Status and dependencies\\
\midrule\endhead
Hahn normal forms and finite workspace localization &
Standard background. The interpretation in actual surreal numbers uses
Conway normal forms, as discussed in \cite{BG,Repo}.\\[5pt]
Positive-support evaluation &
Classical support principle, deduced here from Higman's word theorem
\cite{Higman}. It does not assert sequential convergence.\\[5pt]
Prony uniqueness, Hankel factorization, Hermite duals &
Classical finite algebra; proved here for completeness and credited to the
Prony/interpolation background \cite{BY,KDB}.\\[5pt]
Theorem~\ref{prony:thm:main} and Proposition~\ref{prony:prop:last} &
Principal proposed refinement: exact arbitrary-rank, pointwise threshold,
nonlinear weight control, and a universal highest-moment witness. Depends
on cofactor linear algebra, simple lifting, and divided differences.\\[5pt]
Proposition~\ref{prony:prop:rows} &
Exact valuation consequence of the classical Hermite inverse, with explicit
cancellation invariant. Not a claim to have invented the inverse matrix.\\[5pt]
Proposition~\ref{prony:prop:lattice} &
Support-based arbitrary-rank polynomial adaptation of the established
differential precision principle \cite{CRV}.\\[5pt]
Corollary~\ref{prony:cor:nonlinear} &
Exact precision-set application, under its extra smallness criterion.
Not asserted to hold throughout the entire sharp-threshold region.\\[5pt]
Theorem~\ref{prony:thm:graph} &
Proposed sufficient nonuniform certificate. It uses cycle potentials and
strongly summable matrix inversion. Its conditions are not claimed
necessary.\\[5pt]
Verification program &
Exact checks of finite examples only. Not a formal proof, a proof of
infinite support assertions, or a priority certificate.\\[5pt]
Section~\ref{prony:subsec:collection} &
Added on placement in the collection: provenance, notation, the meaning of
``moment'', and cross-references to neighbouring reports. No mathematical
claim.\\[5pt]
Remark~\ref{prony:rem:coefficient} &
Added on placement: an elementary comparison of $\Theta$ with the threshold
$T(P)$ of \lbl{polynomial:thm:stability} \cite{RepoPoly}, with two small
examples. No new theorem, and no coefficient-side universal sharpness
claimed.\\
\bottomrule
\end{longtable}

\section{Reproducibility and source audit}
\label{prony:app:audit}
The repository was inspected through its connected GitHub interface. The
pinned commit is recorded in Section~\ref{prony:subsec:repository}. At that
commit the inspected
\lbl{docs/README.md} listed twenty-six principal reports; the
\lbl{docs/surcomplex/polynomial-algebra/README.md} described its
sixty-one-page merged report, including its coefficient threshold,
Hermite interpolation, and support results. The full main-report collection,
the eighteen separate ZIPs then in \lbl{docs/new}, and removed historical
sources were not exhaustively compared. Indexed no-match results are
reported as observations, not as proof that the terms or results occur
nowhere in the repository. These are statements about the pinned snapshot.
The collection has since grown, this report being one of the additions, and
the archives then in \lbl{docs/new} have been placed into reports
(Section~\ref{prony:subsec:repository}).

The primary literature comparison used the authors' arXiv manuscripts
\cite{BY,KDB,AGY,CRV,BG} and publisher metadata for the journal references.
The exact arbitrary-rank theorem stated here was not identified in that
material. Keyword searches combining Prony with Hahn, surreal, or valued
fields did not yield a matching primary theorem, but several search results
were broad or irrelevant; they cannot establish comprehensive absence.
The manuscript therefore does not certify priority. The manuscript modified
no repository file; joining the collection added only this report's own
directory.

\begin{samepage}
Compile from this directory with
\begin{verbatim}
latexmk -pdf -halt-on-error -interaction=nonstopmode article.tex
python code/verify.py
\end{verbatim}
\end{samepage}
Run the second command on a copy of the directory: it deletes and rewrites
\lbl{data/verification.txt} (Section~\ref{prony:sec:verification}). The
tested SymPy version is pinned in \lbl{data/requirements.txt}.
The LaTeX bibliography is internal. No proprietary fonts, external images,
or downloaded papers are required or bundled. The verification script uses
SymPy; the PDF and its proofs do not require executing that script.

\begin{thebibliography}{9}
\bibitem{BY}
D.~Batenkov and Y.~Yomdin,
\emph{On the accuracy of solving confluent Prony systems},
SIAM Journal on Applied Mathematics \textbf{73}(1) (2013), 134--154.
\href{https://doi.org/10.1137/110836584}{doi:10.1137/110836584}.
Author manuscript: \href{https://arxiv.org/abs/1106.1137}{arXiv:1106.1137}.

\bibitem{KDB}
R.~Katz, N.~Diab, and D.~Batenkov,
\emph{On the accuracy of Prony's method for recovery of exponential sums
with closely spaced exponents},
Applied and Computational Harmonic Analysis \textbf{73} (2024), 101687.
\href{https://doi.org/10.1016/j.acha.2024.101687}{doi:10.1016/j.acha.2024.101687}.
Author manuscript: \href{https://arxiv.org/abs/2302.05883}{arXiv:2302.05883},
version 2, 16 April 2024.


\bibitem{AGY}
A.~Akinshin, G.~Goldman, and Y.~Yomdin,
\emph{Geometry of error amplification in solving Prony system with
near-colliding nodes},
\href{https://arxiv.org/abs/1701.04058}{arXiv:1701.04058}, version 5,
17 December 2019. See Definition 1.3 and the discussion of the Prony curve
$S_{2d-2}$.

\bibitem{CRV}
X.~Caruso, D.~Roe, and T.~Vaccon,
\emph{Tracking $p$-adic precision},
LMS Journal of Computation and Mathematics \textbf{17}, special issue A
(2014), 274--294.
Author manuscript: \href{https://arxiv.org/abs/1402.7142v1}{arXiv:1402.7142v1},
28 February 2014. See Lemma~3.4 and Proposition~3.12: these concern
complete fields with real-valued non-Archimedean norms and Banach spaces.
The arbitrary-rank result here is proved separately by Hahn support
evaluation, rather than an application of that normed-space theorem.

\bibitem{BG}
O.~Bournez and Q.~Guilmant,
\emph{Surreal fields stable under exponential and logarithmic functions},
2022 preprint, \href{https://arxiv.org/abs/2201.08199}{arXiv:2201.08199}.
Used here for background on surreal normal forms and Hahn subfields,
not for any Prony reconstruction assertion.

\bibitem{Higman}
G.~Higman,
\emph{Ordering by divisibility in abstract algebras},
Proceedings of the London Mathematical Society (3) \textbf{2}(1)
(1952), 326--336.
\href{https://doi.org/10.1112/plms/s3-2.1.326}{doi:10.1112/plms/s3-2.1.326}.

\bibitem{Repo}
V.~Reshetnikov and contributors,
\emph{Surreal: Lean formalization and research reports}, GitHub repository,
inspected 22 September 2026, commit
\texttt{4cf691c7d951e037739d32d9f5c387dcce724f3c}.
\href{https://github.com/VladimirReshetnikov/Surreal/tree/4cf691c7d951e037739d32d9f5c387dcce724f3c}{Pinned repository snapshot}.
Specifically the documentation catalogue, formalization-ledger excerpts,
and polynomial-algebra report README. Research reports are not treated
as independently refereed literature.

\bibitem{RepoPoly}
\emph{Surcomplex Polynomial Algebra: Root Geometry, Newton Profiles,
Multiscale Stability, and Residue Duality}, report
\lbl{docs/surcomplex/polynomial-algebra} of this collection; in particular
\lbl{polynomial:thm:stability}, \lbl{polynomial:ex:sharpdisc},
\lbl{polynomial:cor:discprecision} and \lbl{polynomial:thm:crt}.
Cited on placement; an AI-assisted research draft, not refereed literature.

\bibitem{RepoSpec}
\emph{Surcomplex Spectral Theory}, report
\lbl{docs/surcomplex/spectral-theory} of this collection; its section on
Hankel square-class profiles. Cited on placement; an AI-assisted research
draft.

\bibitem{RepoMeas}
Report \lbl{docs/surreal/hahn-valued-measures-and-probability} of this
collection, on Hahn-valued measures, moment representation and probability.
Cited on placement; an AI-assisted research draft.

\bibitem{RepoHerg}
Report \lbl{docs/surcomplex/hahn-herglotz-positivity} of this collection,
on Hahn--Herglotz normalization and trigonometric moments. Cited on
placement; an AI-assisted research draft.

\bibitem{RepoAna}
\emph{Surcomplex Analysis}, report \lbl{docs/surcomplex/analysis} of this
collection; in particular \lbl{d:thm:moments}. Cited on placement; an
AI-assisted research draft.
\end{thebibliography}
\end{document}
